% META: {"id": "TCOMPL-INEG-CO-008", "chapitre": "TCOMPL-INEGALITES", "type_objet": "corrige", "exercice_ref": "TCOMPL-INEG-EX-008", "capacites_codes": ["C3"], "sources_inspiration": [], "status": "generated"}
% BEGIN-VERIFY
% from sympy import Rational
% points = [(Rational(0), Rational(0)), (Rational(1, 2), Rational(1, 5)),
%           (Rational(1), Rational(1))]
% aire = sum(Rational(points[i][1] + points[i + 1][1], 2)
%            * (points[i + 1][0] - points[i][0]) for i in range(2))
% assert aire == Rational(35, 100)
% gini = 2 * (Rational(1, 2) - aire)
% assert gini == Rational(3, 10)
% assert 2 * (Rational(1, 2) - Rational(1, 2)) == 0
% assert 2 * (Rational(1, 2) - 0) == 1
% END-VERIFY
\begin{corrige}{TCOMPL-INEG-EX-008}
\textbf{1.} L'aire sous la ligne brisée est la somme de deux trapèzes :
$\tfrac{0+0{,}2}{2}\times0{,}5+\tfrac{0{,}2+1}{2}\times0{,}5
=0{,}05+0{,}3=0{,}35$.

\textbf{2.} L'aire entre la bissectrice et la courbe vaut
$0{,}5-0{,}35=0{,}15$, donc $G=2\times0{,}15=0{,}3$.

\textbf{3.} Égalité parfaite : la courbe est la bissectrice, l'aire entre
elles est nulle et $G=0$. Concentration totale : la courbe suit l'axe
horizontal puis remonte verticalement, l'aire vaut $\tfrac12$ et $G=1$.
L'indice de Gini varie donc entre $0$ et $1$.
\end{corrige}
